\chapter{Layer 8: Optimization, Search, and the P/NP Swarm}
\label{ch:pnp}

\section{Purpose and Contract}
\label{sec:pnp-purpose}

Layer~8 does two things: it \emph{optimizes} representations and it
\emph{searches} for proofs. Its inbound contract is ``an obligation and a
resource budget''; its outbound contract is ``a witness (a solution or a
better form) or a recorded failure.'' Two sibling modules implement it:
\texttt{prism-skills} (canonical sealing of solved skills) and
\texttt{pnp-attack} (a multi-agent proof-search swarm).

\section{PRISM Skills: Sealed Capabilities}
\label{sec:pnp-prism}

The \texttt{layers/prism-skills/} module's contract is:
\begin{quote}
\textbf{Layer 8: PRISM Skills} \\
Purpose: Canonical JSON, SHA-256d, and WORM sealing. \\
Status: Verified.
\end{quote}
A \emph{skill} here is a solved sub-problem: a transformation that is known to
work, sealed so it can be reused without re-proving. Each skill is encoded as
canonical JSON, double-hashed with SHA-256d, and written to the WORM ledger.
The neighboring \texttt{../math-skills/} repository carries six such skills
(enumeration, isomorphism, symmetry, Hadamard, probabilistic, circuits), each
with a master seal --- a concrete instance of this layer in production.

\section{The P/NP Swarm}
\label{sec:pnp-swarm}

The \texttt{layers/pnp-attack/} module realizes the central swarm insight:
\begin{quote}
\itshape Finding a solution is NP-hard. Verifying a solution is P-time. The
repository only accepts P-verifiable proofs.
\end{quote}
Agents claim problems, compute witnesses (the hard NP work), submit them with
proofs, and the repository verifies (the easy P work) before promoting a
solution. A claim ledger is append-only; a solution pool holds candidate
witnesses; a verifier promotes the first verified one. The constellation's
\texttt{agentos/} substrate (Chapter~\ref{ch:constellation}) is the runtime
that schedules these agents.

\section{The Malice Layer}
\label{sec:pnp-malice}

Adversarial robustness is first-class. \texttt{malice\_layer2/} is a
refutation module: it attempts to \emph{break} claims, modeling an opponent
(M4LKIT in the constellation lore) probing the Bifrost layer. Its
\texttt{metadata.json} carries \texttt{family\_id: 9} and a
\texttt{corpora\_path} of \texttt{novel\_methods/malice\_refutation/},
indicating it is the dedicated red-team boundary. A system that only proves
itself correct is not safe; one that is subjected to active refutation is.

\section{Optimization}
\label{sec:pnp-opt}

Beyond search, Layer~8 optimizes the canonical form for the execution layer:
selecting the cheapest evaluation order, hoisting invariants, and choosing
closed forms from the pattern registry (\cref{ch:lean}) when available. The
optimization is itself a solved skill when repeatable, feeding back into
\texttt{prism-skills}.

\section{Why a Swarm and Not a Solver}
\label{sec:pnp-why}

A single automated theorem prover is a point of failure and a performance
bottleneck. A swarm parallelizes the NP-hard search across agents, each
competing on a slice, while the repository's P-time verifier remains the sole
arbiter of truth. This matches the architecture's recurring pattern: hard work
is distributed, checking is centralized and cheap.

\section{A Worked P/NP Solve}
\label{sec:pnp-worked}

Consider the open problem of certifying the closed form for
$P(k)=k^4$. A swarm agent claims it via the claim ledger and submits a
Lean proof term. The verifier runs:
\begin{enumerate}
  \item \texttt{lean} on the submitted proof; if it type-checks, the
        theorem is admitted to the corpus.
  \item a recomputation of the WORM receipt hash for the new artifact.
  \item a check that the new receipt signs under the Plasma Gate.
\end{enumerate}
On success the problem moves to \texttt{solved}, the convergence log
records a \texttt{problem\_solved} event, and the universe sum advances.
The reward (per the registry) may be a new skill unlock, which is then
loadable by later agents via the inverted-skills memory
(\cref{ch:artifact}).
